\chapter{Recomendaciones}

A partir de la planificación y el inicio del desarrollo del Sistema de Gestión Integral
Universitario (que incluye los módulos de Gestión de Evidencias SINAES y Gestión de
Tiempos de Jornada), y con el objetivo de asegurar su éxito a largo plazo y su evolución
óptima, se presentan las siguientes recomendaciones estratégicas:

\subsection{Enfoque Técnico y Arquitectónico}
\begin{itemize}
    \item \textbf{Priorizar el diseño modular y escalable de la arquitectura} para
        ambos módulos desde las fases iniciales. Esto permitirá la fácil integración
        de nuevas funcionalidades y futuras expansiones sin requerir rediseños
        mayores. Se recomienda enfocarse en:
        \begin{itemize}
            \item Definición clara de interfaces y APIs entre componentes.
            \item Uso de microservicios o arquitecturas distribuidas para el módulo de Tiempos de Jornada.
            \item Estrategias de desacoplamiento para minimizar dependencias.
            \item Adopción de patrones de diseño que faciliten la escalabilidad horizontal.
            \item Consideración de un diseño "API-first" para futuras integraciones.
        \end{itemize}

    \item \textbf{Establecer una estrategia robusta de respaldo y recuperación de datos}
        desde el inicio del proyecto. Esto es crucial para la información de evidencias
        (Google Drive) y los datos de tiempos de jornada (MongoDB). La estrategia debe incluir:
        \begin{itemize}
            \item Definición de la frecuencia de respaldos (diarios, semanales).
            \item Mecanismos de almacenamiento de respaldos en ubicaciones seguras y geográficamente dispersas.
            \item Protocolos de cifrado para datos en reposo y en tránsito.
            \item Pruebas periódicas y documentadas de restauración de datos.
            \item Un plan de recuperación ante desastres (DRP) detallado.
        \end{itemize}

    \item \textbf{Planificar la integración con la API de Google Drive} con un enfoque
        en la robustez y manejo de excepciones. Dado que el módulo de evidencias no
        tendrá un backend propio en la primera fase, es vital asegurar:
        \begin{itemize}
            \item Manejo de límites de cuota y errores de la API.
            \item Estrategias de reintento con retroceso exponencial.
            \item Monitoreo de la salud de la conexión y las operaciones con Drive.
            \item Auditoría de cambios y sincronización para prevenir inconsistencias.
            \item Documentación exhaustiva de los flujos de interacción con Drive.
        \end{itemize}

    \item \textbf{Investigar y planificar futuras capacidades móviles} para ambos módulos,
        considerando la creciente necesidad de acceso desde dispositivos. Esto podría incluir:
        \begin{itemize}
            \item Diseño responsive para el frontend web.
            \item Exploración de frameworks para desarrollo de aplicaciones híbridas.
            \item Identificación de funcionalidades clave para un acceso móvil prioritario.
            \item Consideración de modos offline para la captura de datos en áreas sin conectividad.
            \item Análisis de las necesidades de rendimiento y experiencia de usuario móvil.
        \end{itemize}
\end{itemize}

\subsection{Estrategias Operativas y de Adopción}
\begin{itemize}
    \item \textbf{Diseñar un programa de capacitación integral} para los usuarios
        finales (Gestores de Calidad, Decanos, Directores de Carrera) desde las etapas
        tempranas del desarrollo. Este programa debe:
        \begin{itemize}
            \item Adaptar los contenidos a los roles específicos y niveles de competencia tecnológica.
            \item Incluir sesiones prácticas y materiales de apoyo (manuales, videos).
            \item Fomentar la participación activa y la resolución de dudas en vivo.
            \item Establecer un cronograma de capacitaciones pre-lanzamiento y post-lanzamiento.
            \item Designar usuarios clave como "campeones" del sistema para soporte inicial.
        \end{itemize}

    \item \textbf{Establecer mecanismos claros de retroalimentación de usuarios}
        y gestión de solicitudes de mejora continua. Esto facilitará la identificación
        temprana de problemas y oportunidades de optimización. Se sugiere:
        \begin{itemize}
            \item Implementar un canal centralizado para reportar errores y sugerencias.
            \item Definir un proceso para el análisis y priorización de estas solicitudes.
            \item Comunicar proactivamente el estado de las mejoras implementadas a los usuarios.
            \item Realizar encuestas de satisfacción y entrevistas con usuarios periódicamente.
            \item Fomentar una cultura de mejora continua en la comunidad de usuarios.
        \end{itemize}

    \item \textbf{Definir Indicadores Clave de Desempeño (KPIs)} para ambos módulos
        desde la fase de análisis. Estos KPIs permitirán medir el impacto del sistema
        en la eficiencia operativa y el cumplimiento de objetivos. Deben incluir:
        \begin{itemize}
            \item Métricas de tiempo (ej. tiempo de clasificación de evidencia, tiempo de registro de jornada).
            \item Métricas de calidad de datos (ej. porcentaje de documentos con metadatos completos).
            \item Métricas de cumplimiento (ej. porcentaje de criterios SINAES con evidencias asociadas).
            \item Métricas de uso del sistema y satisfacción del usuario.
            \item Establecimiento de líneas base y metas realistas.
        \end{itemize}

    \item \textbf{Establecer un comité de gobernanza del proyecto} con representación
        de los principales stakeholders (académicos, administrativos, TI). Este comité
        debería:
        \begin{itemize}
            \item Reunirse periódicamente para supervisar el progreso y la alineación estratégica.
            \item Aprobar cambios significativos en el alcance y prioridades.
            \item Gestionar riesgos de alto nivel y facilitar la resolución de conflictos.
            \item Promover la adopción del sistema a nivel institucional.
            \item Asegurar la disponibilidad de recursos necesarios para el proyecto.
        \end{itemize}
\end{itemize}

\subsection{Crecimiento y Alcance Futuro}
\begin{itemize}
    \item \textbf{Diseñar el sistema con una visión de integración con otros sistemas
        universitarios} a futuro. Aunque no sea parte de la fase inicial, identificar
        puntos de integración potenciales (ej. sistemas de recursos humanos, gestión
        académica, finanzas) es crucial para:
        \begin{itemize}
            \item Evitar silos de información y duplicidad de datos.
            \item Optimizar procesos transversales en la UNA.
            \item Maximizar el valor de la información institucional.
            \item Sentar las bases para una visión de campus inteligente.
            \item Reducir la carga administrativa al automatizar flujos de datos.
        \end{itemize}

    \item \textbf{Explorar el potencial de análisis avanzado y predicción} para
        ambos módulos una vez que se acumulen suficientes datos históricos. Esto
        podría incluir:
        \begin{itemize}
            \item Para evidencias SINAES: Predecir posibles deficiencias en criterios.
            \item Para Tiempos de Jornada: Optimización de la asignación de recursos docentes.
            \item Identificación de patrones y tendencias para la toma de decisiones estratégicas.
            \item Integración con herramientas de Business Intelligence.
            \item Desarrollo de dashboards interactivos con proyecciones.
        \end{itemize}

    \item \textbf{Planificar la expansión de las capacidades de reportería} para
        satisfacer necesidades analíticas avanzadas más allá de los reportes iniciales
        en PDF/Excel. Las mejoras futuras podrían incluir:
        \begin{itemize}
            \item Herramientas de reportes personalizables por el usuario final.
            \item Implementación de visualizaciones de datos dinámicas.
            \item Capacidades de análisis multidimensional (OLAP).
            \item Integración con plataformas de datos abiertos de la UNA.
            \item Generación programada y distribución automática de informes clave.
        \end{itemize}
\end{itemize}

\subsection{Gestión del Conocimiento y Documentación}
\begin{itemize}
    \item \textbf{Mantener una documentación técnica y de usuario actualizada}
        a lo largo de todo el ciclo de vida del proyecto. Esta debe ser una tarea
        continua, no solo al final, para asegurar que refleje la evolución del sistema.
        Incluir:
        \begin{itemize}
            \item Un proceso formal de actualización vinculado a cada sprint y liberación.
            \item Diagramas de arquitectura y flujos de trabajo detallados.
            \item Documentación de decisiones técnicas clave y su justificación.
            \item Ejemplos prácticos y casos de uso para los usuarios.
            \item Un sistema de control de versiones para la documentación.
        \end{itemize}

    \item \textbf{Establecer procedimientos de transferencia de conocimiento}
        efectivos dentro del equipo de desarrollo y hacia el personal de TI de la UNA.
        Esto es crucial para la sostenibilidad y mantenimiento futuro. Se debe considerar:
        \begin{itemize}
            \item Sesiones de walkthrough de código y diseño.
            \item Definir protocolos para el traspaso de responsabilidades.
            \item Mantener registros detallados de configuraciones y personalizaciones.
            \item Documentar soluciones a problemas recurrentes y lecciones aprendidas.
            \item Establecer repositorios centralizados de conocimiento técnico.
        \end{itemize}

    \item \textbf{Crear una base de conocimiento de soporte} para los usuarios y
        el equipo de TI. Esta base permitirá la resolución rápida de incidencias
        y reducirá la dependencia del equipo de desarrollo inicial. Debería:
        \begin{itemize}
            \item Organizar soluciones por categorías intuitivas de problemas frecuentes.
            \item Incluir capturas de pantalla y videos instructivos.
            \item Proporcionar procedimientos paso a paso para la resolución de incidencias.
            \item Implementar un sistema de búsqueda eficiente por palabras clave.
            \item Permitir la contribución y calificación de soluciones por parte de los usuarios.
        \end{itemize}

    \item \textbf{Documentar las lecciones aprendidas} en cada fase del proyecto,
        especialmente en los momentos clave (ej. al finalizar cada sprint, al cierre
        del proyecto). Esta práctica contribuirá al mejoramiento de futuros proyectos
        de desarrollo en la UNA. La documentación debe:
        \begin{itemize}
            \item Analizar aspectos exitosos y áreas de mejora en el proceso.
            \item Identificar factores críticos de éxito replicables.
            \item Documentar obstáculos encontrados y estrategias efectivas de superación.
            \item Evaluar la efectividad de diferentes enfoques de gestión del cambio.
            \item Proporcionar recomendaciones concretas para futuros proyectos tecnológicos.
        \end{itemize}
\end{itemize}

\subsection{Sostenibilidad a Largo Plazo}
Para garantizar la viabilidad y relevancia del sistema en el largo plazo, se
recomienda:

\begin{itemize}
    \item \textbf{Asegurar la asignación de un presupuesto anual dedicado} al
        mantenimiento y mejora continua del sistema una vez puesto en producción.
        Esto debe incluir recursos para actualizaciones tecnológicas, capacitación
        continua y desarrollo de nuevas funcionalidades.

    \item \textbf{Fomentar la capacitación interna del personal de TI de la UNA}
        en las tecnologías utilizadas (React, Next.js, TypeScript, NestJS, MongoDB,
        Prisma) para reducir la dependencia de proveedores externos o del equipo
        de desarrollo inicial, asegurando la capacidad de soporte y evolución interna.

    \item \textbf{Establecer alianzas estratégicas con la academia o centros de
        investigación} de la misma UNA para proyectos de investigación aplicada
        que potencien el sistema a futuro, como la integración de IA para análisis
        predictivos o la automatización de procesos.

    \item \textbf{Evaluar periódicamente la alineación del sistema} con los
        objetivos estratégicos de la UNA, las normativas SINAES y las mejores prácticas
        en gestión universitaria, asegurando su relevancia y pertinencia continua.

    \item \textbf{Promover la participación activa en comunidades de desarrolladores
        y usuarios} de las tecnologías empleadas. Esto facilitará el intercambio
        de conocimientos, la identificación de nuevas funcionalidades y la resolución
        colaborativa de desafíos técnicos.
\end{itemize}